problem:  
Let \((M,g_0)\) be a compact \(n\)-dimensional Riemannian spin manifold admitting a **real Killing spinor** \(\psi_0\) with constant \(\lambda \neq 0\), so that \(g_0\) is Einstein and its metric cone  
\[
(C(M), \bar g_0 = dr^2 + r^2 g_0)
\]
is **Ricci-flat** with a **parallel spinor**.  Suppose we fix a weighted Hölder space \(C^{k,\alpha}_\delta\) on symmetric 2-tensors over \(C(M)\), where \(\delta \in (2-n,0)\), and let 
\[
\Delta_L^{\bar g_0} h = \nabla^*\nabla h + 2\,\mathrm{Rm}_{\bar g_0}*h
\]
denote the **Lichnerowicz Laplacian** governing infinitesimal Ricci-flat deformations.

We call a deformation \((\bar g_\epsilon = \bar g_0 + \epsilon\, h + O(\epsilon^2))\) **AC Ricci-flat** if it remains Ricci-flat up to first order, \(h \in C^{k,\alpha}_\delta\), and obeys the gauge condition \(\nabla^{\bar g_0}\!\!\cdot h = 0\).  Its **asymptotic link metric** is then \(g_\epsilon = g_0 + \epsilon\, h_\infty + O(\epsilon^2)\), where \(h_\infty\) is the homogeneous term in the large-\(r\) expansion of \(h\).

Now define the **boundary linearization**  
\[
\beta_* : \ker(\Delta_L^{\bar g_0}) \cap C^{k,\alpha}_\delta \;\longrightarrow\; \mathcal{E}_{\mathrm{KS}}(M,g_0),
\]
where \(\mathcal{E}_{\mathrm{KS}}(M,g_0)\) consists of infinitesimal Einstein deformations of \(g_0\) that preserve the Killing spinor equation,
\[
D_{g_0}\psi_0 = \lambda \psi_0, \quad \text{up to first-order variation.}
\]

**Question:**  
Is \(\beta_*\) surjective?  In other words, does every infinitesimal Einstein deformation \(h_\infty \in \mathcal{E}_{\mathrm{KS}}(M,g_0)\) arise as \(h_\infty = \mathrm{Asymp}_\infty(h)\) for some decaying AC Ricci-flat deformation \(h\) of the cone metric \(\bar g_0\)?  
If not, can one describe an explicit **obstruction operator**
\[
\mathcal{O}: \mathcal{E}_{\mathrm{KS}}(M,g_0)\longrightarrow \mathrm{coker}(\Delta_L^{\bar g_0})_\delta
\]
whose kernel characterizes the extendable Killing-spinor-preserving deformations?  

Furthermore, in a specific holonomy case (e.g. nearly Kähler 6-manifold \(M\) with Calabi–Yau cone, or nearly parallel \(G_2\) link with Spin(7) cone), does \(\mathcal{O}\) vanish identically or exhibit a holonomy-dependent structure?

Why is it a "good" problem:  
This problem formulates, with analytic precision, the open question of **whether intrinsic Killing-spinor-preserving Einstein deformations on a compact manifold extend to genuine asymptotically conical Ricci-flat deformations on its special-holonomy cone**. It connects **spin geometry, elliptic PDEs, and special-holonomy deformation theory** through the natural boundary map between intrinsic and extrinsic moduli spaces—a mechanism central but not yet understood even in canonical examples like Calabi–Yau and Spin(7) cones.  

The investigation demands detailed analysis of the Lichnerowicz operator’s spectral behavior on conical backgrounds and its interplay with the Dirac operator governing the Killing spinor condition. Success would yield a bridge between **Killing spinor geometry** on the base and **Ricci-flat asymptotics** on the cone, illuminating how spinorial symmetries constrain or enable the existence of desingularized special-holonomy metrics.  

Its statement is simple, geometric, and abstract enough to be widely applicable, yet its resolution involves deep microlocal, analytic, and topological tools, meeting the criteria of a truly “good” research problem.